\section{Por qué se necesita Offline RL}

\begin{importantbox}{Definición de Offline RL}
Dado un \textbf{conjunto de datos estático} $\mathcal{D}$ recolectado por una política de comportamiento desconocida $\pi_\beta$, entrenar, \textbf{sin recolectar datos nuevos}, una política $\pi_\theta$ que maximice la recompensa.
\end{importantbox}

El valor de Offline RL:
\begin{itemize}
    \item Aprovechar los datos ya recolectados por humanos o por sistemas existentes.
    \item La recolección de datos en línea puede ser peligrosa o costosa.
    \item Reutilizar los datos de experimentos, proyectos o robots anteriores.
\end{itemize}

\begin{knowledgebox}{Online vs. Offline RL}
Online RL alterna entre «recolectar datos» y «actualizar la política». Offline RL solo cuenta con un conjunto de datos estático, y el despliegue ocurre hasta que termina el entrenamiento. También es posible entrenar primero offline y después online (modo híbrido).
\end{knowledgebox}

\subsection{Resumen de la sección}

Offline RL resuelve el problema de «cómo aprender una buena política a partir de datos ya existentes, sin interactuar con el entorno».

